\chapter{Conclusiones}
\label{ch:conclusiones}

Tras el análisis, diseño e implementación del Ejercicio Práctico N° 2, se concluye que se han cumplido con éxito los
objetivos propuestos. La configuración del periférico convertidor analógico-digital (ADC1) en un entorno Bare Metal ha
permitido interactuar de manera directa con las señales analógicas del microcontrolador y procesar sus valores
digitales equivalentes.

El ejercicio demuestra la viabilidad de implementar sistemas de adquisición híbridos, donde la conversión de datos se
realiza mediante muestreo activo (polling del bit \texttt{EOC}) para garantizar lecturas continuas e instantáneas,
mientras que la conmutación de los canales de lectura se resuelve de manera asíncrona a través de una interrupción
externa desencadenada por hardware (línea EXTI15). Esta combinación optimiza el tiempo de CPU al delegar la detección
de eventos de control al controlador de interrupciones (NVIC).

Por otro lado, la utilización del sensor de temperatura interno resalta la importancia de los periféricos embebidos
para el diagnóstico y seguridad del chip. El procesamiento de sus valores mediante aproximaciones matemáticas basadas
en aritmética entera demuestra ser una práctica de optimización crucial en firmware para sistemas embebidos de recursos
limitados, evitando la inclusión de costosas bibliotecas de punto flotante.

Finalmente, el diseño de la interfaz visual mediante una escala en la barra de 5 LEDs externos permite comprobar de
forma empírica y en tiempo real las variaciones de temperatura y el comportamiento lineal del potenciómetro de entrada.
Este desarrollo proporciona los fundamentos de adquisición analógica necesarios para abordar implementaciones más
complejas en etapas posteriores de este trabajo práctico.
